\section{แบบฝึกหัดเพิ่มเติมสำหรับเตรียมสอบปลายภาค}
\subsection{รายวิชา 5571102 คณิตศาสตร์วิศวกรรมไฟฟ้า}
\subsubsection{บทที่ 4 อนุกรมและทฤษฎีอนุกรมฟูเรียร์}

\par\noindent\rule{\textwidth}{0.4pt}\par

\section{ชุดฝึกที่ 1}

\subsection{ข้อ 1 ลำดับและพจน์ทั่วไป}

กำหนดลำดับ

\[
6,\;10,\;14,\;18,\ldots
\]

จง

\begin{enumerate}
    \item ระบุชนิดของลำดับและหาผลต่างร่วม \(d\)
    \item หาพจน์ทั่วไป \(a_n\)
    \item หาพจน์ที่ 30 คือ \(a_{30}\)
    \item ตรวจสอบคำตอบโดยแทน \(n=1,2,3\) ลงในพจน์ทั่วไป
\end{enumerate}

\textbf{ทักษะที่วัด:} การจำแนกลำดับและสร้างพจน์ทั่วไป

\par\noindent\rule{\textwidth}{0.4pt}\par

\subsection{ข้อ 2 ผลบวกอนุกรมเลขคณิต}

จงหาผลบวก

\[
7+12+17+\cdots+102
\]

โดยแสดงขั้นตอนดังต่อไปนี้

\begin{enumerate}
    \item ระบุ \(a_1\) และ \(d\)
    \item หาจำนวนพจน์ \(N\)
    \item หาผลบวก \(S_N\)
    \item แสดงสูตรที่ใช้และแทนค่าอย่างเป็นขั้นตอน
\end{enumerate}

\textbf{ทักษะที่วัด:} การย้อนหาจำนวนพจน์ก่อนคำนวณผลบวก

\par\noindent\rule{\textwidth}{0.4pt}\par

\subsection{ข้อ 3 อนุกรมเรขาคณิตและการลู่เข้า}

พิจารณาอนุกรม

\[
18+9+4.5+2.25+\cdots
\]

จง

\begin{enumerate}
    \item หาอัตราส่วนร่วม \(r\)
    \item พิจารณาว่าอนุกรมลู่เข้าหรือลู่ออก พร้อมให้เหตุผล
    \item ถ้าลู่เข้า จงหาผลบวกอนันต์ \(S_\infty\)
    \item จงหาผลบวก 5 พจน์แรก \(S_5\)
    \item เปรียบเทียบ \(S_5\) กับ \(S_\infty\) และอธิบายว่าผลที่ได้สอดคล้องกับแนวคิดการลู่เข้าอย่างไร
\end{enumerate}

\textbf{ทักษะที่วัด:} ไม่เพียงแทนสูตร แต่ต้องตีความเงื่อนไข \(|r|<1\)

\par\noindent\rule{\textwidth}{0.4pt}\par

\subsection{ข้อ 4 คาบ ความถี่ และฮาร์มอนิกของสัญญาณ}

สัญญาณแรงดันเป็นคาบมี

\[
T=5\ \text{ms}
\]

จงหา

\begin{enumerate}
    \item ความถี่มูลฐาน \(f_0\)
    \item ความถี่เชิงมุมมูลฐาน \(\omega_0\)
    \item ความถี่ของฮาร์มอนิกลำดับที่ 2
    \item ความถี่ของฮาร์มอนิกลำดับที่ 5
    \item ความถี่ของฮาร์มอนิกลำดับที่ 9
    \item อธิบายสั้น ๆ ว่าเหตุใดความถี่ฮาร์มอนิกจึงไม่ใช่ค่าความถี่แบบสุ่ม
\end{enumerate}

\textbf{เน้น:} นักศึกษาต้องแปลง \(5\,\text{ms}\) เป็นวินาทีก่อนคำนวณ

\par\noindent\rule{\textwidth}{0.4pt}\par

\subsection{ข้อ 5 การจำแนกฟังก์ชันคู่–คี่และผลต่ออนุกรมฟูเรียร์}

จงจำแนกฟังก์ชันต่อไปนี้เป็น \textbf{ฟังก์ชันคู่ ฟังก์ชันคี่ หรือไม่เป็นทั้งคู่และคี่} พร้อมแสดงการตรวจสอบจาก \(f(-t)\)

\[
\begin{aligned}
f_1(t)&=t^5\\
f_2(t)&=t^4+3\\
f_3(t)&=2-t\\
f_4(t)&=\cos(4t)+\cos(2t)\\
f_5(t)&=\sin(3t)-\sin(t)
\end{aligned}
\]

สำหรับแต่ละฟังก์ชัน ให้ระบุเพิ่มเติมว่า จากสมมาตรสามารถสรุปเกี่ยวกับ

\[
a_0,\qquad a_n,\qquad b_n
\]

ได้อย่างไร

\par\noindent\rule{\textwidth}{0.4pt}\par

\section{ชุดฝึกที่ 2}

ชุดนี้ให้ระดับใกล้ข้อสอบจริงขึ้นเล็กน้อย แต่ยังอยู่ในกรอบเดียวกับแบบฝึกหัดท้ายบทข้อ 1–5

\subsection{ข้อ 1 ลำดับเลขคณิตแบบลดลง}

กำหนด

\[
15,\;12,\;9,\;6,\ldots
\]

จง

\begin{enumerate}
    \item หา \(d\)
    \item หา \(a_n\)
    \item หา \(a_{20}\)
    \item หาค่า \(n\) ที่ทำให้ \(a_n=0\) หากมี
\end{enumerate}

\par\noindent\rule{\textwidth}{0.4pt}\par

\subsection{ข้อ 2 ผลบวกอนุกรมเลขคณิตแบบลดลง}

จงหาผลบวก

\[
120+113+106+\cdots+1
\]

โดยต้องแสดงวิธีหาจำนวนพจน์ก่อนหาผลบวก

\par\noindent\rule{\textwidth}{0.4pt}\par

\subsection{ข้อ 3 อนุกรมเรขาคณิตสลับเครื่องหมาย}

กำหนด

\[
20-10+5-2.5+\cdots
\]

จง

\begin{enumerate}
    \item หา \(r\)
    \item ตรวจสอบการลู่เข้า
    \item หาผลบวกอนันต์
    \item หาผลบวก 6 พจน์แรก
    \item อธิบายว่าเหตุใดผลบวกย่อยจึงสามารถสลับสูง–ต่ำ แต่ยังลู่เข้าสู่ค่าหนึ่งได้
\end{enumerate}

\par\noindent\rule{\textwidth}{0.4pt}\par

\subsection{ข้อ 4 วิเคราะห์ฮาร์มอนิกจากคาบ}

สัญญาณกระแสหนึ่งมีคาบ

\[
T=8\ \text{ms}
\]

จงหา

\[
f_0,\qquad \omega_0,\qquad f_3,\qquad f_4,\qquad f_{10}
\]

จากนั้นตอบว่า ถ้าพบองค์ประกอบที่ความถี่ \(500\text{ Hz}\) องค์ประกอบดังกล่าวเป็นฮาร์มอนิกลำดับที่เท่าใด

\par\noindent\rule{\textwidth}{0.4pt}\par

\subsection{ข้อ 5 วิเคราะห์สมมาตรก่อนคำนวณฟูเรียร์}

พิจารณาฟังก์ชัน

\[
\begin{aligned}
f_1(t)&=|t|\\
f_2(t)&=t^7-t\\
f_3(t)&=3+\cos(2t)\\
f_4(t)&=\sin(4t)+\sin(2t)\\
f_5(t)&=1+t^2+\sin(t)
\end{aligned}
\]

จง

\begin{enumerate}
    \item จำแนกแต่ละฟังก์ชันว่าเป็นคู่ คี่ หรือไม่เป็นทั้งคู่และคี่
    \item ระบุสัมประสิทธิ์ที่สามารถสรุปได้ทันทีว่าเป็นศูนย์
    \item ระบุว่าฟังก์ชันใดคาดว่าจะเหลือเฉพาะพจน์โคไซน์
    \item ระบุว่าฟังก์ชันใดคาดว่าจะเหลือเฉพาะพจน์ไซน์
    \item อธิบายว่าฟังก์ชันที่ไม่เป็นทั้งคู่และคี่ เหตุใดจึงไม่สามารถตัด \(a_n\) หรือ \(b_n\) ทั้งชุดทิ้งได้จากสมมาตรเพียงอย่างเดียว
\end{enumerate}

\par\noindent\rule{\textwidth}{0.4pt}\par

\subsection{หมายเหตุสำหรับการเตรียมสอบ}

แบบฝึกหัดทั้ง 2 ชุดออกแบบให้ยึดแนวเดียวกับแบบฝึกหัดท้ายบทข้อ 1–5 ของบทที่ 4 โดยเน้นทักษะสำคัญ ได้แก่

\begin{itemize}
    \item ลำดับและพจน์ทั่วไป
    \item ผลบวกของอนุกรมเลขคณิต
    \item อนุกรมเรขาคณิตและการลู่เข้า
    \item คาบ ความถี่มูลฐาน และฮาร์มอนิก
    \item สมมาตรของฟังก์ชันคู่–คี่และผลต่อสัมประสิทธิ์ฟูเรียร์
\end{itemize}
